% ============================================================================
%  教师版讲义样式（整册小册子）
% ----------------------------------------------------------------------------
%  用法：主文件 03-讲义/讲义-教师版.tex 已经写好，一般不用动。
%        每课一个文件放在 03-讲义/教师版/ 下，开头写
%            \jhlesson{课次}{课题}{本课信息}
%        主文件里用 \input 引进来。
%  编译：XeLaTeX，编译两遍（目录需要第二遍才正确）
% ----------------------------------------------------------------------------
%  设计原则：务实，少加不必要的框。
%    只保留 4 个方框：mubiao 本课目标 / definition 定义 / example 例 / xiaojie 本课小结
%    其余一律排成普通段落：caozuo 讲台操作 / zhu 注 / jielun 结论
% ============================================================================

% ---------------- 字体（跨平台自动选）----------------
\IfFontExistsTF{SimSun}{
  \setCJKmainfont{SimSun}[BoldFont=SimHei,ItalicFont=KaiTi,BoldItalicFont=SimHei]
  \setCJKsansfont{Microsoft YaHei}
  \setCJKmonofont{FangSong}
  \setCJKfamilyfont{zhsong}{SimSun}
  \setCJKfamilyfont{zhhei}{SimHei}
  \setCJKfamilyfont{zhkai}{KaiTi}
  \setCJKfamilyfont{zhfs}{FangSong}
}{
  \IfFontExistsTF{Songti SC}{
    \setCJKmainfont{Songti SC}[BoldFont=Heiti SC,ItalicFont=Kaiti SC,BoldItalicFont=Heiti SC]
    \setCJKsansfont{PingFang SC}
    \setCJKmonofont{STFangsong}
    \setCJKfamilyfont{zhsong}{Songti SC}
    \setCJKfamilyfont{zhhei}{Heiti SC}
    \setCJKfamilyfont{zhkai}{Kaiti SC}
    \setCJKfamilyfont{zhfs}{STFangsong}
  }{
    \setCJKmainfont{FandolSong-Regular.otf}[BoldFont=FandolHei-Regular.otf,ItalicFont=FandolKai-Regular.otf,BoldItalicFont=FandolHei-Regular.otf]
    \setCJKsansfont{FandolHei-Regular.otf}
    \setCJKmonofont{FandolFang-Regular.otf}
    \setCJKfamilyfont{zhsong}{FandolSong-Regular.otf}
    \setCJKfamilyfont{zhhei}{FandolHei-Regular.otf}
    \setCJKfamilyfont{zhkai}{FandolKai-Regular.otf}
    \setCJKfamilyfont{zhfs}{FandolFang-Regular.otf}
  }
}
\providecommand{\songti}{}\renewcommand{\songti}{\CJKfamily{zhsong}}
\providecommand{\heiti}{}\renewcommand{\heiti}{\CJKfamily{zhhei}}
\providecommand{\kaishu}{}\renewcommand{\kaishu}{\CJKfamily{zhkai}}
\providecommand{\fangsong}{}\renewcommand{\fangsong}{\CJKfamily{zhfs}}
\IfFontExistsTF{TeX Gyre Pagella}{%
  \setmainfont{TeX Gyre Pagella}%
}{\setmainfont{Latin Modern Roman}}

% ---------------- 宏包 ----------------
\usepackage[a4paper,left=2.3cm,right=2.3cm,top=2.5cm,bottom=2.2cm,headheight=15pt]{geometry}
\usepackage{amsmath,amssymb,mathtools,bm}
\usepackage{booktabs,array,multirow}
\PassOptionsToPackage{dvipsnames,svgnames}{xcolor}
\usepackage{graphicx,tikz}
\usepackage{enumitem}
\usepackage{fancyhdr}
\usepackage[many]{tcolorbox}
\tcbuselibrary{breakable,skins}
\usepackage[hidelinks]{hyperref}

% ---------------- 图片搜索路径 ----------------
\graphicspath{{./}{./图片/}{../../00-共享/}{../00-共享/}{00-共享/}{00-共享/图片/}}

% ---------------- 配色（与课件同一族绿色）----------------
\definecolor{jyblue}{RGB}{0,90,160}
\definecolor{jygreen}{RGB}{0,100,0}
\definecolor{jyteal}{RGB}{0,121,107}
\definecolor{jygray}{RGB}{85,85,85}
\definecolor{jydark}{RGB}{20,90,20}

% ---------------- 页眉页脚 ----------------
\newcommand{\jhcourse}{从数数开始的抽象代数生活}
\newcommand{\jhbookname}{教师用书}
\newcommand{\jhauthor}{姜博瀚}
\newcommand{\jhsetbook}[2]{\renewcommand{\jhcourse}{#1}\renewcommand{\jhbookname}{#2}}
\newcommand{\jhcurrentlesson}{}

\pagestyle{fancy}
\fancyhf{}
\fancyhead[L]{\small\songti\jhcourse}
\fancyhead[R]{\small\songti\jhcurrentlesson}
\fancyfoot[C]{\small\thepage}
\renewcommand{\headrulewidth}{0.4pt}
\renewcommand{\footrulewidth}{0pt}
% 章首页/目录页用的 plain 样式：不写页眉页脚（目录页因此不带页码）
\fancypagestyle{plain}{\fancyhf{}\renewcommand{\headrulewidth}{0pt}}

% ---------------- 封面 ----------------
% 主文件里写： \jhcover{教师用书}{给站在讲台上的你}{一句副题}
\newcommand{\jhcover}[3]{%
  \thispagestyle{empty}%
  \begingroup
  \centering
  \vspace*{1.4cm}
  {\color{jydark}\rule{\textwidth}{1.8pt}}\\[5pt]
  {\color{jydark}\rule{\textwidth}{0.6pt}}\\[2.5cm]
  {\heiti\fontsize{31pt}{40pt}\selectfont 从数数开始的抽象代数生活}\\[18pt]
  {\songti\fontsize{13pt}{20pt}\selectfont\color{jygray} 厦门大学附属实验中学\quad 校本课程}\\[2.1cm]
  {\color{jydark}\rule{0.42\textwidth}{0.8pt}}\\[16pt]
  {\heiti\fontsize{23pt}{30pt}\selectfont\color{jydark} #1}\\[10pt]
  {\songti\fontsize{12.5pt}{19pt}\selectfont\color{jygray} #2}\\[14pt]
  {\color{jydark}\rule{0.42\textwidth}{0.8pt}}\\[3.1cm]
  {\songti\fontsize{12pt}{18pt}\selectfont 授课教师：\jhauthor}\\[6pt]
  {\songti\fontsize{11pt}{17pt}\selectfont\color{jygray} \today}\\[1.5cm]
  {\kaishu\fontsize{13pt}{20pt}\selectfont\color{jydark} #3}%
  \vfill
  \par
  \endgroup
}

% ---------------- 目录 ----------------
\renewcommand{\contentsname}{目\hspace{1em}录}

% ---------------- 每课的开头 ----------------
% \jhlesson{课次}{课题}{本课信息（时长、教具、是否发单页等，不写就留空 {}）}
\newcommand{\jhlessonhead}[2]{%
  \begin{center}
    {\color{jydark}\rule{0.92\textwidth}{1.1pt}}\\[13pt]
    {\heiti\fontsize{19pt}{24pt}\selectfont 第 #1 课}\\[9pt]
    {\heiti\fontsize{16pt}{23pt}\selectfont #2}\\[13pt]
    {\color{jydark}\rule{0.92\textwidth}{1.1pt}}
  \end{center}
}
\newcommand{\jhlesson}[3]{%
  \par
  \ifdim\pagetotal>0pt\clearpage\fi
  \phantomsection
  \addcontentsline{toc}{chapter}{第 #1 课\quad #2}%
  \renewcommand{\jhcurrentlesson}{第 #1 课\quad #2}%
  \markboth{\jhcourse}{\jhcurrentlesson}%
  \vspace*{0.4cm}%
  \jhlessonhead{#1}{#2}%
  \vspace{-4pt}%
  \ifx\relax#3\relax\else
    \begin{center}{\small\songti\color{jygray} #3}\end{center}%
  \fi
  \vspace{6pt}%
}

% ---------------- 四个方框 ----------------
\tcbset{
  jybase/.style={
    enhanced, breakable, sharp corners,
    boxsep=3pt, left=9pt, right=8pt, top=6pt, bottom=6pt,
    before skip=9pt, after skip=9pt,
    fonttitle=\heiti,
    attach title to upper={\ },
  }
}
\newcommand{\jysub}[1]{\ifx\relax#1\relax\else（#1）\fi}

\newtcolorbox{mubiao}[1][]{jybase,
  colback=jyblue!6!white, colframe=jyblue!6!white, coltitle=jyblue!60!black,
  title={本课目标\jysub{#1}},
  overlay={\draw[line width=2.6pt, jyblue] (frame.north west) -- (frame.south west);}}
\newtcolorbox{definition}[1][]{jybase,
  colback=jyblue!6!white, colframe=jyblue!6!white, coltitle=jyblue!60!black,
  title={定义\jysub{#1}},
  overlay={\draw[line width=2.6pt, jyblue] (frame.north west) -- (frame.south west);}}
\newtcolorbox{example}[1][]{jybase,
  colback=jyteal!6!white, colframe=jyteal!6!white, coltitle=jyteal!60!black,
  title={例\jysub{#1}},
  overlay={\draw[line width=2.6pt, jyteal] (frame.north west) -- (frame.south west);}}
\newtcolorbox{xiaojie}[1][]{jybase,
  colback=jydark!7!white, colframe=jydark!7!white, coltitle=jydark!60!black,
  title={本课小结\jysub{#1}},
  overlay={\draw[line width=2.6pt, jydark] (frame.north west) -- (frame.south west);}}

% ---------------- 不加框的三类段落 ----------------
% 讲台操作、注、结论：只留文本本身，不加框、不上色
\newenvironment{caozuo}[1][]{\par\vspace{1pt}\ignorespaces}{\par\vspace{1pt}}
\newenvironment{zhu}[1][]{\par\vspace{1pt}\ignorespaces}{\par\vspace{1pt}}
\newenvironment{jielun}[1][]{\par\vspace{1pt}\ignorespaces}{\par\vspace{1pt}}

% ---------------- 图片占位符 ----------------
\newcommand{\imgph}[2][4cm]{%
  \begin{center}%
  \begin{tikzpicture}%
    \node[draw=jydark, densely dashed, line width=1pt, rounded corners=4pt,
          fill=jydark!5!white, inner sep=8pt, align=center, text width=0.72\textwidth,
          minimum height=#1, font=\small]
      {\textcolor{jydark}{\bfseries 【图片占位】}\\[4pt] #2};%
  \end{tikzpicture}%
  \end{center}%
}

% ---------------- 列表与行距 ----------------
\setlist[itemize]{itemsep=2pt,topsep=3pt,parsep=0pt,leftmargin=1.6em}
\setlist[enumerate]{itemsep=2pt,topsep=3pt,parsep=0pt,leftmargin=2.0em}
\linespread{1.15}